% !TEX root = ../frode-bibelen.tex
% ----------------------------------------------------------------------------
%  Frode-bibelen -- 9 Filer
%  Inkluderes fra frode-bibelen.tex (kompiler alltid hovedfilen).
% ----------------------------------------------------------------------------

% ========================================================== 9 FILER ========
\kapittel{Filer}

\begin{regel}[{Gullregel: filformat}]
Les \textbf{nøyaktig} det formatet oppgaven beskriver -- og skriv \textbf{nøyaktig} det samme når du skriver til fil.
\cc{xxxLesFraFil} og \cc{xxxSkrivTilFil} skal være speilbilder av hverandre.
\end{regel}

\begin{center}
\begin{tabular}{@{}L{5.6cm}L{10.0cm}@{}}
\toprule
\textbf{På fila står \dots} & \textbf{Les med \dots} \\
\midrule
ett tall & \cc{fscanf(inn, "\%i", \&x);} \\
flere tall/tegn på samme linje & \cc{fscanf(inn, "\%i \%c \%f", \&(k->tlf), \&tegn, \&(k->saldo));} \\
blank før tekst / \cc{'\bs n'} etter tall & \cc{getc(inn);} \ (leser og forkaster \emph{ett} tegn) \\
tekst ut linja (kan ha blanke) & \cc{fgets(buffer, STRLEN, inn);} \cc{buffer[strlen(buffer)-1] = '\bs 0';} \\
ett ord (ingen blanke) & \cc{fscanf(inn, "\%s", buffer);} \\
\cc{0}/\cc{1} som skal bli \cc{bool} & \cc{fscanf(inn, "\%i", \&kan);} \cc{k->dag[i] = (kan == 1);} \\
bokstav som skal bli \cc{enum} & \cc{fscanf(inn, " \%c", \&tegn);} \cc{k->type = (tegn == 'P') ? Privat : Bedrift;} \\
antall + så mange verdier & \cc{fscanf} antallet, deretter en \cc{for}-løkke med \cc{fscanf}, til slutt \cc{getc}. \\
\bottomrule
\end{tabular}
\end{center}

\subsection{Les fra fil -- antallet står først}\label{sec:lesfil}
Filformat (eksempel):
\begin{kode}[language={},backgroundcolor=\color{white}]
3
91234567 P 250.00 Ola Hansen
45678901 B 0.00 Fjellsport AS
99887766 P 99.50 Kari Nordmann
\end{kode}
\begin{kode}
void lesFraFil()  {
  FILE* innfil = fopen("KUNDER.DTA", "r");             //  Åpner for lesing.
  if (innfil)  {                                       //  Filen finnes:
     fscanf(innfil, "%i", &gAntallKunder);             //  Leser antallet,
     getc(innfil);                                     //    og '\n' etter det.
     for (int i = 0;  i < gAntallKunder;  i++)  {      //  Lager og leser alle:
         gKunder[i] = (struct Kunde*) malloc(sizeof(struct Kunde));
         kundeLesFraFil(innfil, gKunder[i]);           //  Leser selv.
     }
     fclose(innfil);                                   //  Lukker INNI if-en.
     printf("\n\tKundene er lest inn fra 'KUNDER.DTA'.\n");
  } else
     printf("\n\tFant ikke filen 'KUNDER.DTA'!\n");
}

void kundeLesFraFil(FILE* inn, struct Kunde* kunde)  {
  char buffer[STRLEN];
  char tegn;
  fscanf(inn, "%i %c %f", &(kunde->tlf), &tegn, &(kunde->saldo));
  getc(inn);                                           //  Forkaster ' ' før navnet.
  fgets(buffer, STRLEN, inn);                          //  Resten av linja = navn.
  buffer[strlen(buffer)-1] = '\0';                     //  Fjerner '\n'.
  kunde->navn = (char*) malloc((strlen(buffer)+1) * sizeof(char));
  strcpy(kunde->navn, buffer);
  kunde->type = (tegn == 'P') ? Privat : Bedrift;      //  Bokstav -> enum.
}
\end{kode}

\subsection{Les fra fil -- ukjent antall (les til slutten)}
Filformat: navn på én linje, resten på neste. Les \textbf{første post før løkka} og \textbf{neste post sist i løkka} (H23).
\begin{kode}
void lesFraFil2()  {
  char buffer[STRLEN];
  FILE* innfil = fopen("KUNDER.DTA", "r");
  if (innfil)  {
     gAntallKunder = 0;
     fgets(buffer, STRLEN, innfil);                    //  Prøver å lese 1. navn.
     while (!feof(innfil)  &&  gAntallKunder < MAXKUNDER)  {
       buffer[strlen(buffer)-1] = '\0';
       gKunder[gAntallKunder] = (struct Kunde*) malloc(sizeof(struct Kunde));
       kundeLesFraFil2(innfil, gKunder[gAntallKunder], buffer);
       gAntallKunder++;
       fgets(buffer, STRLEN, innfil);                  //  Prøver å lese neste navn.
     }
     fclose(innfil);
  } else
     printf("\n\tFant ikke filen 'KUNDER.DTA'!\n");
}

void kundeLesFraFil2(FILE* inn, struct Kunde* kunde, const char* navn)  {
  char tegn;
  kunde->navn = (char*) malloc((strlen(navn)+1) * sizeof(char));
  strcpy(kunde->navn, navn);
  fscanf(inn, "%i %c %f", &(kunde->tlf), &tegn, &(kunde->saldo));
  getc(inn);                                           //  Forkaster '\n'.
  kunde->type = (tegn == 'P') ? Privat : Bedrift;
}
\end{kode}

\subsection{Skriv til fil}\label{sec:skrivfil}
\begin{kode}
void skrivTilFil()  {
  FILE* utfil = fopen("KUNDER.DTA", "w");              //  "w" OVERSKRIVER filen!
  if (utfil)  {
     fprintf(utfil, "%i\n", gAntallKunder);            //  Antallet først,
     for (int i = 0;  i < gAntallKunder;  i++)         //    så hver kunde:
         kundeSkrivTilFil(utfil, gKunder[i]);
     fclose(utfil);
  } else
     printf("\n\tKlarte ikke aa aapne 'KUNDER.DTA'!\n");
}

void kundeSkrivTilFil(FILE* ut, const struct Kunde* kunde)  {
  fprintf(ut, "%i %c %.2f %s\n", kunde->tlf,           //  SAMME format som lesingen.
          (kunde->type == Privat) ? 'P' : 'B', kunde->saldo, kunde->navn);
}
\end{kode}

\begin{center}
\begin{tabular}{@{}L{4.6cm}L{11.0cm}@{}}
\toprule
\textbf{Funksjon} & \textbf{Merk} \\
\midrule
\cc{fopen(navn, "r" | "w" | "a")} & Les / skriv (tømmer filen først) / legg til på slutten. \cc{NULL} hvis det ikke gikk. \\
\cc{fclose(fil)} & Alltid når du er ferdig -- men bare hvis \cc{fopen} lyktes. \\
\cc{fprintf(fil, format, ...)} & Som \cc{printf}, men til fil. \\
\cc{fscanf(fil, format, ...)} & Som \cc{scanf}. Hopper over blanke og \cc{'\bs n'} \emph{før} tall og \cc{\%s}. \\
\cc{fgets(buf, len, fil)} & Leser til og med \cc{'\bs n'} (eller \cc{len}-1 tegn). \textbf{Beholder} \cc{'\bs n'}. \\
\cc{getc(fil)} & Leser ett tegn (brukes for å «spise» blank eller \cc{'\bs n'}). \\
\cc{feof(fil)} & Sann først \textbf{etter} et lesforsøk som traff slutten. \\
\bottomrule
\end{tabular}
\end{center}

\begin{felle}[{Fil-feller}]
\begin{itemize}
  \item \cc{fscanf} et tall og så \cc{fgets}, \textbf{uten \cc{getc}} imellom: \cc{fgets} leser bare den tomme resten av linja (\cc{"\bs n"}).
  \item Glemt å fjerne \cc{'\bs n'} etter \cc{fgets} -- da skrives navnet med linjeskift og \cc{strcmp} finner ingenting.
  \item \cc{while (!feof(fil))} \emph{uten} å lese før løkka -- siste post leses dobbelt.
  \item \cc{fclose} utenfor \cc{if (fil)} -- \cc{fclose(NULL)} kan krasje.
  \item Test skrivingen med et annet filnavn (Frode brukte f.eks.\ \cc{.DT2}) så du ikke ødelegger datafila.
  \item Bruk filnavnet \textbf{nøyaktig} slik oppgaven oppgir det.
\end{itemize}
\end{felle}
